% 07_cirugia.tex
\begin{landscape}

	% =============================================================================
	% CONFIGURACIÓN
	% =============================================================================
	\setlength{\tabcolsep}{2pt}
	\setlength{\extrarowheight}{3pt}
	\fontsize{10pt}{10pt}\selectfont

	% =============================================================================
	% CÁLCULO DE ANCHOS DE COLUMNAS
	% =============================================================================
	\setlength{\availablewidth}{\linewidth}

	% Restamos (8 columnas * 2 lados) = 16 espacios para un ajuste matemático exacto
	\addtolength{\availablewidth}{-16\tabcolsep}

	% Restamos el ancho de las líneas verticales (9 líneas * 0.4pt estándar)
	\addtolength{\availablewidth}{-3.6pt}

	\newcommand{\tableheader}{\rowcolor{blue!10}
		\textbf{ESPECIE} & \textbf{Nº DE CASOS} & \textbf{CIRUGÍA} &
		\textbf{FÁRMACO} & \textbf{PRINCIPIO ACTIVO} & \textbf{DOSIS} &
		\textbf{VÍA DE ADMINISTRACIÓN} & \textbf{OBSERVACIO\-NES} \\
	}

	% =============================================================================
	% TABLA
	% =============================================================================
	\begin{longtable}{
		|>{\centering\arraybackslash}m{0.08\availablewidth}    %Especie
		|>{\centering\arraybackslash}m{0.07\availablewidth}    %Casos
		|>{\centering\arraybackslash}m{0.10\availablewidth}    %Cirugia
		|>{\raggedright\arraybackslash}m{0.16\availablewidth}  %Farmaco
		|>{\raggedright\arraybackslash}m{0.22\availablewidth}  %Principio
		|>{\raggedright\arraybackslash}m{0.13\availablewidth}  %Dosis
		|>{\raggedright\arraybackslash}m{0.12\availablewidth}  %Via
		|>{\raggedright\arraybackslash}m{0.12\availablewidth}| %Observaciones
	}

		% -----------------------------------------------------------------------------
		% CAPTION Y LABEL
		% -----------------------------------------------------------------------------
		\caption{\raggedright Registro de Cirugías realizadas durante el periodo de internado}
		\label{tab:cirugias} \\
		\hline
		\tableheader
		\hline
		\endfirsthead

		% -----------------------------------------------------------------------------
		% ENCABEZADO (PRIMERA PÁGINA Y SIGUIENTES)
		% -----------------------------------------------------------------------------
		\multicolumn{8}{l}{\normalsize\textbf{(Continuación) Cuadro \ref{tab:cirugias}}. Registro de Cirugías realizadas durante el periodo de internado} \\[6pt]
		\hline
		\tableheader
		\hline
		\endhead

		% -----------------------------------------------------------------------------
		% PIE DE CONTINUACIÓN
		% -----------------------------------------------------------------------------
		\hline
		\multicolumn{8}{r}{\footnotesize\textit{Continúa en la siguiente página \ldots}}
		\endfoot

		\hline
		\multicolumn{8}{l}{\footnotesize\textit{\textbf{Fuente}: Elaboración propia.}}
		\endlastfoot

		% =============================================================================
		% CANINO – ORQUIECTOMÍA
		% =============================================================================
		\multirow{9}{*}{\textbf{Canino}} &
		\multirow{9}{*}{3} &
		\multirow{9}{*}{Orquiectomía} &
		Xilacina & Xilacina & 1\,ml/20\,kg & IV-IM & \\
		\cline{4-7}
		& & & Lidocaína & Lidocaína & Según criterio médico & IM-SC & \\
		\cline{4-7}
		& & & Keta-A & Ketamina 10\% & 1\,ml/10\,kg & IV-IM & \\
		\cline{4-7}
		& & & Cefavet & Cefalexina & 1\,ml/20\,kg & IM-SC & \\
		\cline{4-7}
		& & & Galmetrim plus & Cipermetrina, imidacloprid, sulfadiazina de plata & Según criterio médico & Local externo & \\
		\cline{4-7}
		& & & Agua oxigenada & Peróxido de hidrógeno & Según criterio médico & Tópico & \\
		\cline{4-7}
		& & & Yodo podovidona & Complejo de yodo molecular con podovidona & Según criterio médico & Tópico & \\
		\cline{4-7}
		& & & Ácido poliglicólico 3/0 & Ácido poliglicólico & --- & --- & \\
		\cline{4-7}
		& & & Hilo externo & Lino retorcido & --- & --- & \\
		\hline

		% =============================================================================
		% CANINO – OVH
		% =============================================================================
		\multirow{17}{*}{\textbf{Canino}} &
		\multirow{17}{*}{7} &
		\multirow{17}{*}{OVH} &
		Xilacina & Xilacina & 1\,ml/20\,kg & IV-IM & \\
		\cline{4-7}
		& & & Lidocaína & Lidocaína & Según criterio médico & IM-SC & \\
		\cline{4-7}
		& & & Keta-A & Ketamina 10\% & 1\,ml/10\,kg & IV-IM & \\
		\cline{4-7}
		& & & Cefavet & Cefalexina & 1\,ml/20\,kg & IM-SC & \\
		\cline{4-7}
		& & & Galmetrim plus & Cipermetrina, imidacloprid, sulfadiazina de plata & Según criterio médico & Local externo & \\
		\cline{4-7}
		& & & Agua oxigenada & Peróxido de hidrógeno & Según criterio médico & Tópico & \\
		\cline{4-7}
		& & & Yodo podovidona & Complejo de yodo molecular con podovidona & Según criterio médico & Tópico & \\
		\cline{4-7}
		& & & Desalgina & Dipirona & 1\,ml/10\,kg & IV-IM-SC & \\
		\cline{4-7}
		& & & Atropina & Sulfato de atropina & 0,5\,ml/10\,kg & IV-IM-SC & \\
		\cline{4-7}
		& & & Antihem & Vitamina K & 2-5\,ml & IV & \\
		\cline{4-7}
		& & & Galerfin & Clorfeniramina maleato & 0,2-2\,ml & IV-IM & \\
		\cline{4-7}
		& & & Mercepton & Acetil D-L-Metionina, cloruro de colina, Inositol, Vitamina B1 B2 B6 B12, nicotinamida, pantotenato de calcio & 1\,ml/10\,kg & IV-IM-SC & \\
		\cline{4-7}
		& & & Ringer Lactato & Electrolitos & Según criterio médico & IV & \\
		\cline{4-7}
		& & & Ácido poliglicólico 2/0 & Ácido poliglicólico & --- & --- & \\
		\cline{4-7}
		& & & Ácido poliglicólico 3/0 & Ácido poliglicólico & --- & --- & \\
		\cline{4-7}
		& & & Hilo externo & Lino retorcido & --- & --- & \\
		\hline

		% =============================================================================
		% FELINO – ORQUIECTOMÍA
		% =============================================================================
		\multirow{9}{*}{\textbf{Felino}} &
		\multirow{9}{*}{9} &
		\multirow{9}{*}{Orquiectomía} &
		Xilacina & Xilacina & 1\,ml/20\,kg & IV-IM & \\
		\cline{4-7}
		& & & Lidocaína & Lidocaína & Según criterio médico & IM-SC & \\
		\cline{4-7}
		& & & Keta-A & Ketamina 10\% & 1\,ml/10\,kg & IV-IM & \\
		\cline{4-7}
		& & & Cefavet & Cefalexina & 1\,ml/20\,kg & IM-SC & \\
		\cline{4-7}
		& & & Galmetrim plus & Cipermetrina, imidacloprid, sulfadiazina de plata & Según criterio médico & Local externo & \\
		\cline{4-7}
		& & & Agua oxigenada & Peróxido de hidrógeno & Según criterio médico & Tópico & \\
		\cline{4-7}
		& & & Yodo podovidona & Complejo de yodo molecular con podovidona & Según criterio médico & Tópico & \\
		\cline{4-7}
		& & & Ácido poliglicólico 3/0 & Ácido poliglicólico & --- & --- & \\
		\cline{4-7}
		& & & Hilo externo & Lino retorcido & --- & --- & \\
		\hline

		% =============================================================================
		% FELINO – OVH
		% =============================================================================
		\multirow{17}{*}{\textbf{Felino}} &
		\multirow{17}{*}{14} &
		\multirow{17}{*}{OVH} &
		Xilacina & Xilacina & 1\,ml/20\,kg & IV-IM & \\
		\cline{4-7}
		& & & Lidocaína & Lidocaína & Según criterio médico & IM-SC & \\
		\cline{4-7}
		& & & Keta-A & Ketamina 10\% & 1\,ml/10\,kg & IV-IM & \\
		\cline{4-7}
		& & & Cefavet & Cefalexina & 1\,ml/20\,kg & IM-SC & \\
		\cline{4-7}
		& & & Galmetrim plus & Cipermetrina, imidacloprid, sulfadiazina de plata & Según criterio médico & Local externo & \\
		\cline{4-7}
		& & & Agua oxigenada & Peróxido de hidrógeno & Según criterio médico & Tópico & \\
		\cline{4-7}
		& & & Yodo podovidona & Complejo de yodo molecular con podovidona & Según criterio médico & Tópico & \\
		\cline{4-7}
		& & & Desalgina & Dipirona & 1\,ml/10\,kg & IV-IM-SC & \\
		\cline{4-7}
		& & & Atropina & Sulfato de atropina & 0,5\,ml/10\,kg & IV-IM-SC & \\
		\cline{4-7}
		& & & Antihem & Vitamina K & 2-5\,ml & IV & \\
		\cline{4-7}
		& & & Galerfin & Clorfeniramina maleato & 0,2-2\,ml & IV-IM & \\
		\cline{4-7}
		& & & Mercepton & Acetil D-L-Metionina, cloruro de colina, Inositol, Vitamina B1 B2 B6 B12, nicotinamida, pantotenato de calcio & 1\,ml/10\,kg & IV-IM-SC & \\
		\cline{4-7}
		& & & Ringer Lactato & Electrolitos & Según criterio médico & IV & \\
		\cline{4-7}
		& & & Ácido poliglicólico 2/0 & Ácido poliglicólico & --- & --- & \\
		\cline{4-7}
		& & & Ácido poliglicólico 3/0 & Ácido poliglicólico & --- & --- & \\
		\cline{4-7}
		& & & Hilo externo & Lino retorcido & --- & --- & \\
		\hline

	\end{longtable}
	\vspace{-6pt}\normalsize
	% Resumen conciso de la tabla
	La tabla registra las cirugías realizadas durante el internado en el Hospital para Animales. Se documentaron 33 procedimientos: 3 orquiectomías y 7 OVH en caninos, y 9 orquiectomías y 14 OVH en felinos. Los protocolos quirúrgicos incluyen anestesia con Xilacina, Lidocaína y Keta-A, antibióticos como Cefavet, tratamientos tópicos como Galmetrim plus y Yodo podovidona, suturas con Ácido poliglicólico y Hilo externo. En las OVH, se agregó soporte con analgésicos como Desalgina, Atropina, fluidoterapia con Ringer Lactato, protectores hepáticos como Mercepton y vitaminas, con resultados satisfactorios en todos los casos.

\end{landscape}